\documentclass[../main.tex]{subfiles}
\ifSubfilesClassLoaded{\setcounter{chapter}{9}}{}
\begin{document}

\chapter{Pemrograman Berorientasi Objek (OOP) Dasar}

\begin{subcpmk}
  \item Sub-CPMK 7.1: Mendefinisikan kelas dengan atribut dan metode, membuat objek, serta menerapkan pewarisan sederhana untuk memodelkan entitas program
\end{subcpmk}

\SesiRPS{11}{7.1}{$\sim$170 menit}{CPMK-7}

\section{Sarana dan prasarana}

Komputer laboratorium dengan Python 3.10+, editor/IDE, terminal, dan akses ke dokumentasi daring terbatas sesuai kebijakan prodi.

\section{Keselamatan kerja dan etika}

Tidak ada bahan kimia berisiko; tetap jaga postur ergonomis. Etika kode: tulis program sendiri atau sebutkan sumber jika memakai contoh dari tutorial.

\section{Prosedur kerja}

\begin{enumerate}[leftmargin=*]
  \item Baca ringkasan teori kelas dan objek di bawah ini.
  \item Replikasi contoh di Python, lalu modifikasi nama atribut dan metode.
  \item Implementasikan hierarki sederhana (induk--anak) sesuai aktivitas.
  \item Uji objek dari REPL atau skrip terpisah; simpan berkas sesuai konvensi penamaan.
\end{enumerate}

\section{Konsep Kelas dan Objek}

\textbf{Kelas} adalah cetak biru; \textbf{objek} adalah instansi. Di Python, atribut biasanya diinisialisasi di \texttt{\_\_init\_\_}.

\begin{lstlisting}[caption=Kelas sederhana dengan metode]
class Mahasiswa:
    def __init__(self, nim, nama):
        self.nim = nim
        self.nama = nama

    def perkenalan(self):
        return f"{self.nim} - {self.nama}"

m = Mahasiswa("12345", "Andi")
print(m.perkenalan())
\end{lstlisting}

\section{Pewarisan Sederhana}

Anak mewarisi metode induk dan dapat menimpa (\textit{override}).

\begin{lstlisting}[caption=Pewarisan]
class Pengguna:
    def __init__(self, username):
        self.username = username

class Admin(Pengguna):
    def __init__(self, username, level):
        super().__init__(username)
        self.level = level

    def info(self):
        return f"{self.username} (admin level {self.level})"
\end{lstlisting}

\section{Studi Kasus dan Contoh Kode}

Berikut adalah contoh-contoh program Python yang mengimplementasikan konsep-konsep pada modul ini.

\subsection{Kode: \texttt{kendaraan.py}}
\begin{lstlisting}[language=Python,caption={\texttt{kendaraan.py}}]
"""Contoh kelas dasar dan pewarisan."""


class Kendaraan:
    def __init__(self, merk: str, tahun: int) -> None:
        self.merk = merk
        self.tahun = tahun

    def info(self) -> str:
        return f"{self.merk} ({self.tahun})"

class Mobil(Kendaraan):
    def __init__(self, merk: str, tahun: int, jumlah_pintu: int) -> None:
        super().__init__(merk, tahun)
        self.jumlah_pintu = jumlah_pintu

    def info(self) -> str:
        dasar = super().info()
        return f"{dasar} - {self.jumlah_pintu} pintu"

def main() -> None:
    v1 = Kendaraan("Yamaha", 2015)
    v2 = Mobil("Toyota", 2022, 4)
    print(v1.info())
    print(v2.info())

if __name__ == "__main__":
    main()
\end{lstlisting}

\section{Referensi sesi}

\begin{itemize}[leftmargin=*]
  \item \textbf{[13]} Downey, A. B. \textit{Think Python, 2e} --- bab berorientasi objek.
  \item \textbf{[19]} Real Python --- tutorial kelas dan OOP Python 3.
  \item \textbf{[8, 9]} \textit{The Python Tutorial} (docs.python.org) --- kelas.
\end{itemize}

\begin{aktivitas}
  \item Buat kelas \texttt{Produk} dengan atribut \texttt{nama}, \texttt{harga}, dan metode \texttt{ringkasan()} yang mengembalikan string deskriptif.
  \item Buat kelas \texttt{ProdukDiskon} yang mewarisi \texttt{Produk} dan\newline menambah atribut \texttt{persen\_diskon}; timpa \texttt{ringkasan()} agar menyebut harga setelah diskon.
  \item Uji dengan minimal dua objek di skrip \texttt{NIM\_TIF105\_M11\_oop.py}.
\end{aktivitas}

\begin{latihan}
  \item Jelaskan perbedaan kelas dan objek dalam satu paragraf.
  \item Kapan \texttt{super()} berguna dalam pewarisan?
  \item Apa konvensi \texttt{self} pada metode instance?
\end{latihan}

\begin{asesmen}
\textbf{Formatif (checklist dosen/asisten):} kelas berjalan tanpa error; \texttt{\_\_init\_\_} mengikat atribut; ada minimal satu pewarisan dengan metode tambahan atau override.

\textbf{Kuis singkat:} tuliskan satu baris yang membuat objek dari kelas buatan Anda dan memanggil satu metode.
\end{asesmen}

\begin{checklist}
  \item Saya dapat mendefinisikan kelas dengan \texttt{\_\_init\_\_} dan metode instance
  \item Saya dapat membuat subclass dan memanggil \texttt{super()}
  \item Saya mengumpulkan skrip sesuai format penamaan
\end{checklist}

\begin{rangkuman}
Bab ini memperkenalkan OOP dasar: kelas, objek, \texttt{\_\_init\_\_}, metode, dan pewarisan sederhana dengan \texttt{super()}.

\textbf{Kata kunci:} \code{class}, \code{self}, \code{super()}, pewarisan
\end{rangkuman}

\end{document}
